\documentclass[11pt,a4paper]{article}
\usepackage[utf8]{inputenc}
\usepackage[margin=1in]{geometry}
\usepackage{setspace}
\usepackage{times}

\begin{document}
\begin{center}
    \Large \textbf{Personal Statement} \\
    \vspace{0.2cm}
    \large \textit{University of Michigan-Dearborn}
\end{center}
\vspace{0.5cm}

\setstretch{1.15}

My journey into power engineering did not begin with a fascination for electricity, but with a reluctant compromise. It has been a continuous evolution---from cultivating pure mathematical and programming logic to grasping the physical foundations of power systems, ultimately leading me to find my purpose at the intersection of both disciplines.

As a student at a school for the gifted, I spent my nights relentlessly structuring lines of code for competitive programming tournaments. I was captivated by the absolute control this computational mindset offered. I loved how algorithmic logic could break down immense, chaotic problems into perfectly solvable pieces, allowing me to continuously accelerate processing speeds through creative, boundless optimization. I was certain my future lay entirely in software, but an unforeseen turning point in my life abruptly derailed this plan, necessitating a sudden shift in my trajectory. I chose to pursue an undergraduate degree in Electrical Engineering, a decision that initially felt like a profound compromise---forcing me to trade the limitless freedom of code for a physical discipline strictly bound by rigid, mechanical laws.

In my early university years, studying the foundational laws of power grids felt incredibly constraining. The industry seemed to be nothing more than turning knobs and pulling levers to maintain heavy machinery, leaving my algorithmic mindset entirely sidelined. However, my perspective fundamentally shifted during a Power System Analysis class. Staring at the complex equations, I suddenly realized that behind the physical infrastructure lay a beautiful, intricate order of mathematical logic and power flows. I recognized the greatness of the legacy built by previous generations, yet I longed to use computational thinking to elevate it rather than just passively operating it. From that defining moment, the grid was no longer a collection of lifeless wires, but a colossal mathematical graph waiting to be optimized.

During my time in a progressive research lab, I eagerly attempted to apply decentralized algorithms to real-world grid operations. I quickly encountered the immense historical inertia of a top-down operational mindset. This centralized legacy is a magnificent feat that has successfully ensured energy security for decades, but when faced with the unpredictability of modern renewable energy, it resorts to blunt curtailment rather than flexible optimization. I realized that my algorithmic solutions were clashing directly with a deeply ingrained philosophy of centralized control. This struggle forced me to step back and acknowledge that crafting a perfect algorithm in a laboratory is meaningless unless one respects the immense physical and historical momentum of the grid it aims to serve. I could not simply use micro-level code to patch a system that required a fundamental architectural evolution. This realization extinguished my satisfaction with being a conventional engineer and ignited a profound desire to discover a completely new operational paradigm.

Standing before a corporate environment that explicitly advised me to abandon programming and focus on fixed operational routines and bureaucratic relationship-building, I faced a defining crossroads. I saw that the influx of renewable energy was an irreversible force, acting like a tightly compressed spring against the current infrastructure. While others were content with maintaining the status quo, I knew that the traditional physical boundaries would inevitably be breached by the sheer volume of this new energy landscape. To simply conform would mean turning my back on the computational skills I had spent years cultivating. I understood that to meaningfully address the upcoming energy transition, I needed much sharper, cross-disciplinary weapons that my current environment could not provide. I realized I had to seek out an ecosystem where decentralized markets and algorithmic innovation were not just academic theories, but the foundation of reality.

I actively sought out an academic environment capable of bridging this technological divide. I soon discovered a significant challenge: very few programs successfully integrate high-level software development with high-voltage physical realities, often keeping computer science and power engineering in separate silos. The University of Michigan-Dearborn, however, uniquely resolves this friction; its progressive curriculum and research ecosystem seamlessly erase the boundaries between algorithms and physical infrastructure, offering the exact cross-disciplinary fusion I need. I am determined to immerse myself in UM-Dearborn to become a global systems architect who writes the code that will redefine resilience and stability for the power grids of the future.

\end{document}
